% !TEX program = xelatex
% ============================================================
%  全国大学生数学建模竞赛 (CUMCM) LaTeX 通用论文模板
%  编译方式: XeLaTeX
% ============================================================
\documentclass[12pt,a4paper]{ctexart}

% ======================== 宏包加载 ========================
\usepackage[top=2.5cm, bottom=2.5cm, left=2.8cm, right=2.8cm]{geometry}
\usepackage{amsmath, amssymb, amsthm}
\usepackage{graphicx}
\usepackage{float}
\usepackage{subcaption}
\usepackage{booktabs}
\usepackage{array}
\usepackage{multirow}
\usepackage{longtable}
\usepackage{xcolor}
\usepackage{hyperref}
\hypersetup{
    colorlinks=true,
    linkcolor=blue!70!black,
    citecolor=green!50!black,
    urlcolor=blue!60!black
}
\usepackage{listings}
\lstset{
    basicstyle=\small\ttfamily,
    breaklines=true,
    frame=single,
    numbers=left,
    numberstyle=\tiny\color{gray},
    keywordstyle=\color{blue},
    commentstyle=\color{green!60!black},
    stringstyle=\color{red!70!black},
    tabsize=4
}

\usepackage{setspace}
\onehalfspacing
\usepackage{enumitem}
\setlist{nosep, leftmargin=2em}
\usepackage{titlesec}
\usepackage{caption}
\usepackage{bm}
\usepackage{mathtools}

% ======================== 字体设置 ========================
\setCJKmainfont{SimSun}[BoldFont=SimHei, ItalicFont=KaiTi]
\setCJKsansfont{SimHei}
\setCJKmonofont{FangSong}
\setmainfont{Times New Roman}
% 显式把 ctex 的 \songti/\heiti/\kaishu/\fangsong 族都配上粗体，避免
% 在摘要等使用 \songti 的环境里 \textbf{中文} 无法加粗（显示不出黑体）。
\setCJKfamilyfont{zhsong}{SimSun}[BoldFont=SimHei]
\setCJKfamilyfont{zhhei}{SimHei}[BoldFont=SimHei]
\setCJKfamilyfont{zhkai}{KaiTi}[BoldFont=SimHei]
\setCJKfamilyfont{zhfang}{FangSong}[BoldFont=SimHei]

% ======================== 编号格式 ========================
\renewcommand{\thesection}{\chinese{section}、}
\renewcommand{\thesubsection}{\arabic{section}.\arabic{subsection}}
\renewcommand{\thesubsubsection}{\arabic{section}.\arabic{subsection}.\arabic{subsubsection}}

% ======================== 标题格式 ========================
\titleformat{\section}
  {\centering\heiti\zihao{-3}\bfseries}
  {\thesection}{0em}{}
\titlespacing*{\section}{0pt}{1.5ex plus .5ex minus .2ex}{1.0ex plus .2ex}

\titleformat{\subsection}
  {\heiti\zihao{4}\bfseries}
  {\thesubsection}{1em}{}
\titlespacing*{\subsection}{0pt}{1.2ex plus .4ex minus .2ex}{0.8ex plus .2ex}

\titleformat{\subsubsection}
  {\heiti\zihao{-4}\bfseries}
  {\thesubsubsection}{1em}{}
\titlespacing*{\subsubsection}{0pt}{1.0ex plus .3ex minus .1ex}{0.6ex plus .1ex}

% ======================== 定理环境 ========================
\theoremstyle{definition}
\newtheorem{definition}{定义}[section]
\newtheorem{assumption}{假设}[section]
\theoremstyle{plain}
\newtheorem{theorem}{定理}[section]
\newtheorem{lemma}{引理}[section]
\newtheorem{proposition}{命题}[section]
\newtheorem{corollary}{推论}[section]
\theoremstyle{remark}
\newtheorem{remark}{注}[section]

% ======================== 自定义命令 ========================
\renewenvironment{abstract}
  {\begin{center}\heiti\zihao{-3}摘\quad 要\end{center}\zihao{-4}\songti}
  {\par\vspace{1em}}

\newcommand{\keywords}[1]{%
  \par\noindent\textbf{关键词：}#1}

\newenvironment{symboltable}
  {\begin{center}\zihao{5}\begin{tabular}{>{\centering\arraybackslash}p{3.5cm} p{7.5cm} >{\centering\arraybackslash}p{1.5cm}}
  \toprule
  \textbf{符号} & \textbf{含义} & \textbf{单位} \\
  \midrule}
  {\bottomrule\end{tabular}\end{center}}

\newcommand{\step}[1]{\par\noindent\textbf{Step #1}\quad}

% ============================================================
\begin{document}
\pagestyle{plain}

\begin{center}
{\zihao{-2}\heiti\bfseries 论文标题}
\end{center}
\vspace{1em}

% ======================== 摘要 ========================
\setcounter{page}{1}
\pagenumbering{arabic}

\begin{abstract}
% 第一段：研究背景 + 本文总体思路概述（2~3句）
本文针对……问题，首先建立了……模型，并对……进行了详细分析。

% 之后每段对应一个问题，格式为"针对问题X：方法+关键结果"
\textbf{针对问题一：}本问在给定参数下，建立了……模型，采用……方法（如二分查找法/解析求解）进行求解，最终得到……（关键数值结果）。

\textbf{针对问题二：}在问题一的基础上，本问以……为决策变量，以最大化……为目标建立了优化模型，通过……算法（如多起点变步长搜索/粒子群算法）求解，得到最优策略为……，此时目标函数值为……。此外，本文用……算法进行验证，两种方法的相对误差仅为……，说明算法可靠。

\textbf{针对问题三：}本问将模型拓展至……情形，建立了……优化模型，采用……算法（如差分进化算法）对……进行优化，求解出……（关键结果）。

\textbf{针对问题四：}本问进一步考虑……情形，建立了……模型，提出了……策略（如降维/协同优化），最终求得……。

\textbf{针对问题五：}本问为最复杂的……情形，由于决策变量维度过高，本文采用……策略（如分层次优化/任务分配+参数优化），将问题分解为……，最终求得……。
\end{abstract}

\keywords{关键词1\quad 关键词2\quad 关键词3\quad 关键词4\quad 关键词5}

% ======================== 目录（可选，取消注释启用） ========================
% \newpage
% \tableofcontents
% \newpage

% ============================================================
%  正文
% ============================================================

\section{问题重述}

\subsection{问题背景}

% 写作框架：
% (1) 介绍研究对象/现象的实际背景（是什么、用在哪）；
% (2) 说明该问题的现实意义或应用价值（为什么要研究）；
% (3) 引出本文要解决的核心问题（从背景自然过渡到题目）。
%
% 示例写法：
% ……是……领域中的关键技术，其通过……实现……，具有……的特点，
% 在……方面得到了广泛的应用。然而，如何……仍然是一个亟待解决的
% 问题。本文针对……这一情景，在……条件已经给定的前提下，需要综合
% 考虑……，对……进行优化，以达到……的目标。

\subsection{问题提出}

% 写作框架（参照优秀论文"问题提出"写法）：
% (1) 开头一句话概括已知条件（给了哪些数据/参数），
%     以"本文需要解决以下问题："收尾，引出问题清单；
% (2) 每个问题独立成段，段首加粗"问题X："，
%     每段写清：问题情形 + 需要做什么 + 目标是什么；
% (3) 各问逐层递进：基础计算 → 单因素优化 → 多资源协同
%     → 多主体协同 → 全局优化。

现已知道……的初始参数和……数据，本文需要解决以下问题：

\textbf{问题一：}针对……基础情形，考虑……因素，建立……模型，求解……（基础计算问题）。

\textbf{问题二：}在问题一的基础上，将……推广至一般情形，以最大化……为目标建立优化模型，求解最优策略（单因素优化）。

\textbf{问题三：}考虑……数目增多情形，寻求最优……方案，以延长/提升……（多资源协同）。

\textbf{问题四：}将问题三拓展至多……情形，新增……协同约束，求解最优部署/分配策略（多主体协同）。

\textbf{问题五：}综合多……与多……情形，需同时解决……分配与……调度问题，实现全局最优（最复杂情形）。


\section{问题分析}

\subsection{问题一的分析}

% 分析框架：
% (1) 明确本问要算什么（目标）；
% (2) 梳理已知条件与未知量；
% (3) 说明解题路线：先做什么、再做什么、最后做什么；
% (4) 点明选用该方法的理由。
%
% 示例写法：
% 问题一要求在给定参数下计算……。首先，根据题目条件可以求出……
% 作为后续模型的初值；其次，需要利用……知识建立……的数学关系；
% 最后，基于上述关系，采用……方法（如变步长搜索/解析法）求解……。

\subsection{问题二的分析}

% 分析框架：
% (1) 说明本问与问题一的区别（从"给定参数计算"到"多参数优化"）；
% (2) 明确决策变量是什么、目标函数是什么；
% (3) 分析约束来源（物理约束、边界条件等）；
% (4) 说明为何选择某优化算法（如目标函数无解析表达式→不能用梯度法）。
%
% 示例写法：
% 问题二不同于问题一在给定参数下求解，而是需要将……作为决策变量，
% 以最大化……为目标建立优化模型。由于目标函数没有解析表达式，
% 梯度信息难以获取，因此本文采用……算法进行求解。

\subsection{问题三的分析}

% 分析框架：
% (1) 本问相比前一问新增了什么复杂度（如多弹协同、时序耦合）；
% (2) 分析新复杂度带来的建模变化（判定式如何扩展、约束如何增加）；
% (3) 说明求解策略的调整（如从低维搜索升级为进化算法）。

\subsection{问题四的分析}

% 分析框架：
% (1) 本问与问题三的本质区别（如从单机多弹→多机协同）；
% (2) 新的核心难点（如空间协同、决策变量维度增大）；
% (3) 降维或分解策略的思路来源（如利用几何关系缩小搜索空间）。

\subsection{问题五的分析}

% 分析框架：
% (1) 本问为何不能直接求解（变量维度过高、组合爆炸）；
% (2) 分层/分解策略的核心思想（如先分配任务再优化参数）；
% (3) 各子问题如何独立求解、最终结果如何合并。


\section{模型假设}

\begin{enumerate}[label=\arabic*.]
  \item 假设一。
  \item 假设二。
  \item 假设三。
\end{enumerate}


\section{符号说明}

\begin{symboltable}
$x(t)$ & 示例符号含义 & m \\
$v$ & 示例符号含义 & m/s \\
$\theta$ & 示例符号含义 & rad \\
\end{symboltable}


\section{模型准备}

% 推导后续各问建模所需的公共公式、基本定理、几何关系等。

\subsection{基础公式推导}

\begin{equation}
f(x) = ax^2 + bx + c
\label{eq:example}
\end{equation}

\subsection{核心判定条件}

% 给出贯穿全文的核心数学判据。


\section{模型建立与求解}

\subsection{问题一模型建立及求解}

\subsubsection{模型建立}

% 决策变量、目标函数、约束条件。

\subsubsection{模型求解}

% 算法描述，可用 \step{} 命令：
\step{1} 初始化。

\step{2} 迭代计算。

\step{3} 输出结果。

\subsubsection{结果展示}

% 给出数值结果，配合表格。
\begin{table}[H]
\centering
\caption{表标题}
\label{tab:example}
\zihao{5}
\begin{tabular}{ccc}
\toprule
列1 & 列2 & 列3 \\
\midrule
数据 & 数据 & 数据 \\
\bottomrule
\end{tabular}
\end{table}

\subsubsection{结果分析与验证}

% 敏感性分析、对比验证、合理性讨论。


\subsection{问题二模型建立及求解}

\subsubsection{模型建立}
\subsubsection{模型求解}
\subsubsection{结果展示}
\subsubsection{结果分析与验证}


% 更多问题依此类推：
% \subsection{问题三模型建立及求解}
% \subsection{问题四模型建立及求解}
% ...


\section{模型的评价与推广}

\subsection{模型的优点}

\begin{enumerate}[label=\arabic*.]
  \item 优点一。
  \item 优点二。
\end{enumerate}

\subsection{模型的缺点}

\begin{enumerate}[label=\arabic*.]
  \item 缺点一。
  \item 缺点二。
\end{enumerate}

\subsection{模型的推广}

% 讨论模型在更一般情形下的适用性与改进方向。


% ======================== 参考文献 ========================
\newpage
\begin{thebibliography}{99}
\addcontentsline{toc}{section}{参考文献}

\bibitem{ref1} 作者. 文献标题[J]. 期刊名, 年份, 卷(期).

\bibitem{ref2} 作者. 书名[M]. 出版社, 年份.

\end{thebibliography}


% ======================== 附录 ========================
\newpage
\appendix
% 附录改用"附录 A/B/C"编号，覆盖正文的中文数字编号
\renewcommand{\thesection}{附录 \Alph{section}}
\renewcommand{\thesubsection}{\Alph{section}.\arabic{subsection}}
\renewcommand{\thesubsubsection}{\Alph{section}.\arabic{subsection}.\arabic{subsubsection}}
% 正文编号自带顿号间距为0em，附录编号后需留出间距
\titleformat{\section}
  {\centering\heiti\zihao{-3}\bfseries}
  {\thesection}{1em}{}

\section{支撑文件列表}

\begin{table}[H]
\centering
\zihao{5}
\begin{tabular}{lll}
\toprule
文件名 & 类型 & 说明 \\
\midrule
problem1.py & Python & 问题一求解代码 \\
\bottomrule
\end{tabular}
\end{table}


\section{AI 使用说明}

% 此部分单独成文，见 AI工具使用说明_模板.tex，按需将内容粘贴至此处。


\section{相关代码}

\begin{lstlisting}[language=Python, caption={示例代码}]
# 在此粘贴关键算法代码
\end{lstlisting}


\end{document}
